\documentclass[a4paper,10pt]{article}
\usepackage[francais]{babel}
\usepackage{url}

\selectlanguage{francais}
\input{letter-LSIIT-UdS.sty}

\selectlanguage{francais}
\begin{document}
\name{Stéphane \textsc{Genaud}}
\fonction{Maître de conférences (HDR)}
\telephoneillkirch{+33 (0)3688 54542}
\email{genaud@unistra.fr}


\begin{letter}[Demande de mise à disposition]{à}%
{Alain Beretz, \\Monsieur le Président de l'université de Strasbourg}


Monsieur le Président,


Par la présente, je sollicite votre accord pour être mis à disposition de 
l'\'Ecole Nationale Supérieure d’Informatique pour l'Industrie et l'Entreprise
(ENSIIE), établissement public administratif.

Je souhaite renforcer l'équipe en place sur le campus de Strasbourg.

Je suis persuadé que la présence de cette antenne à Strasbourg est bénéfique
à l'Université de Strasbourg à de multiples égards. En particulier, elle
attire à Strasbourg des étudiants recrutés au niveau national, et dans une
moindre mesure des étudiants étrangers (15 à 20\% des promotions). 
Elle ouvre également une perspective de poursuite d'études locale aux meilleurs
étudiants de nos IUT.
Les étudiants de l'école intègrent pour la plupart, en option en 2$^e$ et 3$^e$ 
année ou en bi-cursus,  des masters de l'UFR Mathématique-informatique ou de la Faculté 
des Sciences Economiques et de Gestion. Ces étudiants sont aussi considérés
par les enseignants-chercheurs comme des éléments intéressants pour un parcours
orienté recherche.

Les interactions entre l'Université de Strasbourg et l'ENSIIE sont donc réelles
et méritent à mon sens qu'on les renforce. J'aimerais m'attacher à cette tâche
et c'est pourquoi je sollicite cette mise à disposition selon les termes de
la convention proposée par l'ENSIIE (document annexe).



Je vous prie d'agréer, Monsieur le Président, mes sentiments respectueux.

\end{letter}
\begin{flushright}
\includegraphics[width=.30\textwidth]{signgenaud.jpg}
\end{flushright}
%\newpage
%\bibliographystyle{plain}dd%\bibliography{julien_blaise_ED}


\end{document}
